\documentclass[a4paper]{article}     
\usepackage{amsfonts}
\usepackage{amsmath}
\usepackage{amssymb}
\usepackage{graphicx}
\usepackage{color}

\newcommand{\Bgsin}{\operatorname{Bgsin}}
\newcommand{\Bgtan}{\operatorname{Bgtan}}
\newcommand{\Bgcos}{\operatorname{Bgcos}}
\newcommand{\Bgcot}{\operatorname{Bgcot}}

\begin{document}

\noindent \large Auteur: Thomas Olszowka \\
\noindent \large Studierichting: Informatica \\
\noindent \large Oefeningenreeks 2B nr. 11 \\



\medskip

\normalsize

$f(x) = \tan\left(\Bgsin\left(\cos\left(\Bgtan\left(\cos\left(\Bgtan\sqrt{3}\right)\right)\right)\right)\right)$ \\

$\Bgtan \sqrt{3} = \dfrac{\pi}{3}$ want $\tan \dfrac{\pi}{3} = \sqrt{3}$ \\

neem $ y = \Bgtan \sqrt{3}$ dan is $\cos y = \cos \dfrac{\pi}{3} = \dfrac{1}{2}$ \\

neem $z = \cos y$ dan is $\Bgtan z = \Bgtan \dfrac{1}{2}$ \\

neem $u = \Bgtan \dfrac{1}{2}$ dan is $\tan u = \dfrac{1}{2}$ \\

we zoeken nu wat $\cos u$ aan gelijk is. \\ 

We kunnen hiervoor de identiteit $\cos^2u = \dfrac{1}{1 + \tan^2 u}$ gebruiken: \\

$\cos^2u = \dfrac{1}{1 + \tan^2 u}$ \\

$\cos u = \sqrt{\dfrac{1}{1 + \tan^2 u}}$ \\

$= \sqrt{\dfrac{1}{1 + \left( \dfrac{1}{2}\right)^2}}$ \\

$= \sqrt{\dfrac{1}{\dfrac{5}{4}}}$ \\

$= \sqrt{\dfrac{4}{5}}$ \\

$= \dfrac{\sqrt{4}}{\sqrt{5}}$ \\

$= \dfrac{2}{\sqrt{5}}$ aangezien $u \in \left]\dfrac{-\pi}{2}, \dfrac{\pi}{2} \right[$ is de cosinus positief. \\

Neem $v = \Bgsin \dfrac{2}{\sqrt{5}}$ dan is $\sin v = \dfrac{2}{\sqrt{5}}$ \\

we zoeken nu wat $\tan v$ aan gelijk is: \\

$\tan v = \dfrac{\sin v}{\cos v}$ \\

hiervoor moeten we nog de cosinus kennen. \\ 

We kunnen de identiteit $\sin^2 v + \cos^2 v = 1$ gebruiken: \\

$\sin^2 v + \cos^2 v = 1$ \\

$\cos^2 v = 1 - \sin^2 v$ \\

$\cos^2 v = 1 - \left(\dfrac{2}{\sqrt{5}}\right)^2$ \\

$= 1 - \dfrac{4}{5}$ \\

$= \dfrac{1}{5}$ \\

$\cos v = \sqrt{\dfrac{1}{5}}$ \\

$= \dfrac{1}{\sqrt{5}}$ aangezien $v \in \left[\dfrac{-\pi}{2}, \dfrac{\pi}{2}\right]$ is de cosinus positief.\\

Nu we de sinus en cosinus hebben gevonden kunnen we de tangens berekenen: \\

$\tan v = \dfrac{\dfrac{2}{\sqrt{5}}}{\dfrac{1}{\sqrt{5}}}$ \\

$= \dfrac{2}{\sqrt{5}} \cdot \dfrac{\sqrt{5}}{1}$ \\

$= \dfrac{2\sqrt{5}}{1\sqrt{5}}$ \\

$= \dfrac{2}{1}$ \\

$= 2$ \\

dus is $\tan\left(\Bgsin\left(\cos\left(\Bgtan\left(\cos\left(\Bgtan\sqrt{3}\right)\right)\right)\right)\right) = 2$ 

\begin{thebibliography}{999}
\end{thebibliography}
\end{document} 